\section{Results}

\subsection{Sample and observed-data modelling population}

The analysis sample was the harmonized baseline (visit V0) of the three FACE
cohorts: $N=9{,}013$ patients (bipolar disorder [BP] $6{,}252$; schizophrenia [SZ]
$2{,}209$; major depression [DR] $552$) recruited across $21$ expert-centre sites
(Table~\ref{tab:characteristics}). The cohorts differ in age, sex distribution and
chronicity, as expected for these disorders. We deliberately fit the measurement
model on the \emph{full} sample without completeness selection. Of the $201$ usable
harmonized common variables (the full role breakdown is in the Methods), the $143$-variable
modelled-indicator pool has mean cell missingness $39.8\%$, preserved as missing and never
imputed; admission to the map is \emph{coverage-agnostic} above a two-cohort floor---the
observed-cell likelihood lets a two-cohort marker serve as an indicator where it is measured,
so about one-third of the fitted indicators are present in only two cohorts.
Supplement~\ref{annA:vars} (Table~\ref{tab:varaccounting}) classifies all $225$ harmonized
variables by role and the selection rule. Diagnosis (cohort and DSM-5 subtype) was held out of the model and used only
as a validation or metadata label. As detailed in the Methods, the FACE-ATLAS map is one global,
cohort-reweighted, full-sample Gaussian-copula fit; because the same structure is recovered
under a plain Gaussian likelihood, the central findings do not depend on that modelling
choice---a robustness result, not an assumption (Supplement~\ref{ann:measurement}). Follow-up visits
at $12$ and $24$ months (retained: $47.4\%$ and $32.8\%$ of V0; \edfig{fig:consort})
were used only to test temporal coherence and prognosis, never to define the map. The
full derivation of the model---the observed-data likelihood, the rotation and identification
argument, the covariance identity and the equivalence with full-information maximum
likelihood---is given in Supplement~\ref{ann:measurement}. Figure~\ref{fig:glance} presents the atlas at a
glance; the study design and the three core modelling principles are detailed in
\edfig{fig:design}.

\begin{figure}[!ht]
  \centering
  \IfFileExists{figures/fig1_glance.png}{%
    \includegraphics[width=0.85\linewidth]{fig1_glance.png}%
  }{%
    \fbox{\parbox[c][6.6cm][c]{0.96\linewidth}{\centering\itshape
      Figure~1 (graphical abstract) renders here once the overview infographic is
      saved to \texttt{article/figures/fig1\_glance.png}.}}%
  }
  \caption{\textbf{Illustration of the FACE ATLAS: latent dimensions of severe mental illness and corresponding archetypes.} The global,
  missingness-aware Bayesian bifactor / exploratory structural-equation model, fit to
  the FACE baseline ($N=9{,}013$) with diagnosis held out, yields an eight-dimension
  transdiagnostic latent space---a general functional-burden factor (\Gfac) and seven
  specific dimensions. The patient cloud is a graded continuum (best silhouette $0.140$,
  indistinguishable from a single-Gaussian null $0.137\pm0.002$), read through five
  archetypal extremes whose convex blends describe every patient: an \emph{activation /
  sleep} corner (A0), a \emph{severe, clean-biology} corner (A1), an
  \emph{immunometabolic} corner (A2; the biology pole), a \emph{trauma / suicidality}
  corner (A3), and a \emph{low-burden} pole (A4). Immunometabolic load is largely
  dissociated from general burden; it is the most durable axis over two years; the
  archetype phenotype improves held-out two-year functional prediction ($\Delta$ELPD
  $+62.8\pm11.2$); and two-year functional remission ranges $22\%$--$63\%$ across
  corners, the immunometabolic corner worst. A1 and A2 carry comparable severity but
  opposite biology.}
  \label{fig:glance}
\end{figure}

\subsection{The fitted object: a posterior coordinate for each patient}

The primary output of FACE-ATLAS is not a cluster label but a \emph{posterior coordinate
distribution}. For patient $i$ the model returns an eight-dimensional coordinate
\begin{equation}
  f_i=\bigl(G_i,\ D_{i,\mathrm{cog}},\ D_{i,\mathrm{immet}},\
  D_{i,\mathrm{sleep}},\ D_{i,\mathrm{mania}},\ D_{i,\mathrm{suic}},\
  D_{i,\mathrm{dev}},\ D_{i,\mathrm{subst}}\bigr)\in\mathbb{R}^{8},
  \label{eq:res-coord}
\end{equation}
one general functional-burden coordinate $G_i$ and seven specific coordinates
(cognition, immunometabolic, sleep, mania, suicidality, developmental risk,
substance use)---the seven specific axes are established as distinct, data-supported
dimensions in \S\ref{subsec:eightdim} (Table~\ref{tab:dimensions}). What the model
actually delivers is the posterior of this coordinate
given the patient's \emph{observed} indicators $x_{i,O_i}$ alone---no imputation, no
completed data matrix (Methods):
\begin{equation}
  p\!\left(f_i\mid x_{i,O_i}\right)\ \approx\ \mathcal{N}\!\left(\hat f_i,\ S_i\right),
  \label{eq:res-post}
\end{equation}
whose mean $\hat f_i$ is the patient's \emph{position} and whose covariance $S_i$ is the
\emph{uncertainty} of that position---a location with error bars, not a category. Two
patients can share an estimated immunometabolic score yet differ in its reliability (a
complete laboratory panel versus a single observed marker), and every downstream analysis
in this paper consumes the full $(\hat f_i,S_i)$ rather than treating coordinates as exact,
so a sparsely measured patient is never mistaken for a precisely located one. The unit of
analysis is thus a position on eight continuous axes with quantified uncertainty, not a
biotype label.

\subsection{How sharply is a patient located? Information accumulation and a minimal-indicator rule}
\label{subsec:localization}

The covariance $S_i$ in \eqref{eq:res-post} is not a fixed property of the model but a
closed-form function of \emph{which} indicators a patient has. Because the indicators are
conditionally independent given the factors, each observed cell contributes additively to the
posterior \emph{precision}: for the continuous (Gaussian-copula) block the posterior covariance
after observing a set of indicators is
\begin{equation}
  S^{-1}=\Phi^{-1}+\sum_{j\ \text{observed}}\frac{1}{\sigma_j^{2}}\,\lambda_j\lambda_j^{\top},
  \qquad
  \hat f=S\!\!\sum_{j\ \text{observed}}\frac{\lambda_j\,x_{ij}}{\sigma_j^{2}},
  \label{eq:res-precision}
\end{equation}
the information form of the EAP score map (Methods, Eq.~\ref{eq:m-scores}), with
$\sigma_j^{2}=1-\lambda_j^{\top}\Phi\lambda_j$ the indicator's residual variance. Two
consequences follow. First, \textbf{the position uncertainty can only shrink as
instruments arrive}: the summed term in \eqref{eq:res-precision} is positive
semidefinite, so every observed cell tightens the ellipsoid, and a patient with no
observed indicators on an axis is returned the prior rather than a false-precise point.
Second, \textbf{the localization is anisotropic}: the rate of collapse on each axis is
set by the loading-weighted information it receives, so a heavily instrumented axis
(immunometabolic, $37$ markers; cognition; general burden) reaches a posterior SD
around $0.2$--$0.3$ from a handful of items, while a thin axis (substance and mania,
two to four indicators) stays wide even at full coverage. This is the quantitative
content of ``a sparsely measured patient is never mistaken for a precisely located
one,'' and Figure~\ref{fig:localization} follows one real, de-identified patient
(BP-62162, bipolar disorder) through the whole pipeline as a single illustrative worked
example. The patient's posterior places them at the extreme of the immunometabolic axis
($+3.57$~SD, posterior SD $0.22$) from only six observed metabolic items, while
cognition and suicidality are prior-dominated (zero items observed, posterior SD
$\approx0.99$; Fig.~\ref{fig:localization}A). The archetype layer summarises this
position as an $86\%$ weight on the immunometabolic corner A2 (Fig.~\ref{fig:localization}B).
Locating the patient in the archetype-by-cohort prognostic atlas
(\S\ref{subsec:eightdim}--\S\ref{subsec:archetypes} below) returns an interval-valued
forecast rather than a point estimate---$27\%$ predicted functional remission within
bipolar disorder, the lowest-remitting archetype in that cohort against $73\%$ at the
low-burden pole (Fig.~\ref{fig:localization}C)---and the model further identifies the
single most informative unmeasured indicator: one cognitive item
(\texttt{cvlt\_short\_delay\_free\_recall}) would raise the cognition axis's reliability
from $0$ to $0.65$ (Fig.~\ref{fig:localization}D). Because this is one patient, the
vignette is illustrative rather than the evidence base; the population-level reliability
curves and cohort-level measurement gaps that generalise it are given next.

\begin{figure}[tbp]
  \centering
  \includegraphics[width=\linewidth]{fig_worked_patient.png}
  \caption{\textbf{A worked patient: coordinates $\to$ archetype $\to$ prognosis $\to$ next
  best measurement.} One real, de-identified patient (BP-62162, bipolar disorder) followed
  end to end through the pipeline; illustrative of the instrument on a single case, not a
  population-level result. \textbf{(A)}~Posterior latent score on all eight axes with $95\%$
  highest-density intervals; the patient sits at the extreme of the immunometabolic axis
  ($+3.57$~SD from six observed items) while cognition and suicidality are prior-dominated
  (no observed items, shown as wide grey intervals). \textbf{(B)}~The archetype layer
  summarises this position as an $86\%$ weight on the immunometabolic corner (A2).
  \textbf{(C)}~Locating the patient in the archetype-by-cohort prognostic atlas gives an
  interval-valued forecast: $27\%$ predicted two-year functional remission within bipolar
  disorder, the lowest of the five archetypes (versus $73\%$ at the low-burden pole A4).
  \textbf{(D)}~Because cognition is currently unmeasured for this patient, the model
  identifies the single most informative next item: one verbal-memory test raises cognition
  reliability from $0$ to $0.65$.}
  \label{fig:localization}
\end{figure}

\paragraph{A loading-dependent minimal-indicator rule.} Equation~\eqref{eq:res-precision} also
answers a practical design question---how many indicators must a factor carry before a patient's
coordinate on it is acceptably precise---and the answer is set not by the raw \emph{count} of
indicators but by the \emph{information} they carry. For a single standardized factor measured by
$m$ conditionally-independent indicators the posterior SD is the accumulated information
\begin{equation}
  \operatorname{SD}(f)=\frac{1}{\sqrt{1+I}},\qquad
  I=\sum_{j=1}^{m}\frac{\lambda_j^{2}}{1-\lambda_j^{2}},
  \label{eq:res-mincount}
\end{equation}
in which each indicator contributes $i(\lambda_j)=\lambda_j^{2}/(1-\lambda_j^{2})$, its
signal-to-noise (reliability) ratio. This per-indicator information is steeply nonlinear in the
loading (Fig.~\ref{fig:mincount}B): $i(0.3)=0.10$, $i(0.5)=0.33$, $i(0.7)=0.96$ and
$i(0.95)=9.3$, so a single strong indicator can be worth dozens of weak ones. Requiring a
posterior SD no larger than a tolerance $\tau$ is the information threshold $I\ge 1/\tau^{2}-1$;
for a target SD~$\le 0.5$ (reliability $0.75$) this is $I\ge 3$, met by about $4$ indicators at
$\lambda=0.7$, $9$ at $\lambda=0.5$ or $31$ at $\lambda=0.3$ for equal loadings
(Fig.~\ref{fig:mincount}A)---or by a single very strong anchor on its own. The cut-off is
therefore loading-dependent, not a fixed instrument count, with minimal count
$m^{\star}=\bigl\lceil(\tfrac{1}{\tau^{2}}-1)(1-\lambda^{2})/\lambda^{2}\bigr\rceil$.

Real factors do not lie on any single equal-loading curve, because their loadings are
heterogeneous, and because the eight axes span four measurement families (Gaussian,
logistic, graded-response, negative-binomial) whose per-indicator information takes a
different closed form in each. Computing the exact Fisher information every indicator
contributes at the population mean under its own family---rather than the
equal-loading approximation above---and adding indicators in decreasing-information
order gives the realised reliability-vs-items curve for each of the eight axes on the
same footing (Fig.~\ref{fig:mincount}, left). Three to six items recover most of each
axis's reliable signal. On the immunometabolic axis, the anthropometric triad---body-mass
index, weight and waist circumference---reaches a score reliability of $0.85$ in three
items and the axis plateaus near $0.88$ (Fig.~\ref{fig:mincount}, right); the remaining
$34$ indicators (the full lipid and inflammatory panel, of the axis's $37$ total markers)
add less than $0.02$ each. Two
axes are intrinsically short: mania/activation, with only two indicators, cannot exceed
a reliability of $0.408$, and substance use, with four, cannot exceed $0.429$. This is a
bank limitation, not a missingness problem---$95$--$96\%$ of patients have at least one
item on each axis---and it is itself an informative finding: those two constructs are
under-instrumented in the current battery, not under-measured in these patients. The
practical lesson for instrument design is therefore not ``collect more questionnaires''
but ``include at least one high-loading marker per intended dimension,'' and, for mania
and substance use specifically, a finer or graded instrument rather than additional
weak items.

\begin{figure}[tbp]
  \centering
  \includegraphics[width=\linewidth]{fig_adaptive_assessment.png}
  \caption{\textbf{Adaptive assessment: how many instruments are enough?} Reliability of
  the axis score against the number of items administered (most-informative first),
  computed from the exact Fisher information each indicator contributes under its own
  likelihood family (Gaussian, logistic, graded-response, negative-binomial), so all
  eight axes are compared on the same footing. \textbf{Left}: all eight axes. Six
  well-instrumented axes climb to reliability $0.85$--$0.90$ within about ten items;
  mania/activation (two indicators) plateaus at $0.408$ and substance use (four
  indicators) at $0.429$---a ceiling set by the size of the instrument bank, not by
  missingness ($95$--$96\%$ of patients have at least one item on each). \textbf{Right}:
  worked example for the immunometabolic axis---the anthropometric triad (BMI, weight,
  waist) reaches reliability $0.85$ in three items and the full $37$-indicator axis
  plateaus near $0.88$.}
  \label{fig:mincount}
\end{figure}

\subsection{A few well-chosen indicators recover most of each dimension}

This reliability-vs-items analysis is, in miniature, an adaptive-assessment result: it
identifies, per dimension, the minimal set of indicators that recovers most of the
reliable signal, so a compact axis-targeted battery can be deployed at a new site with a
stated reliability rather than the full instrument set---directly informing which
measures a future cohort must collect to place patients on the map, and which are
redundant. Because the Fisher information above is evaluated at the population mean, the
curves in Fig.~\ref{fig:mincount} are an upper bound on the efficiency a fully adaptive
test (one that re-evaluates information at each patient's provisional estimate) would
realise; and because the loadings and item parameters are held at their point estimates,
the analysis does not itself propagate parameter uncertainty. Stated proactively, the
mania/activation and substance-use ceilings are a property of the instrument bank the
model has correctly diagnosed, not a failure of the measurement approach: coverage is
not the issue (\S\ref{subsec:localization}), and the fix is a finer or graded severity
instrument, not additional weak binary items.

\paragraph{From uncertainty to a shared battery: designing the next cohort's
measurement.} The same arithmetic that asks ``how many items per axis'' also answers the
consortium-level question: given a fixed measurement budget, which items across
\emph{all} axes should a future cohort collect, and where is each existing cohort
under-measured? By the matrix-determinant lemma applied to \eqref{eq:res-precision},
observing a not-yet-measured indicator $j$ reduces the differential entropy (log-volume)
of the posterior ellipsoid by
$\tfrac12\log\!\big(1+\lambda_j^{\top}S\lambda_j/\sigma_j^{2}\big)$---a closed-form
\emph{value of information} that depends on what has already been measured. Greedily
adding the single most-informative item across all eight axes (recomputed after each,
using the exact per-family Fisher information of \S\ref{subsec:localization}), a
shared $27$-item cross-axis battery brings every axis online at mean reliability $0.70$
(Fig.~\ref{fig:voi}, left); the first item selected on each axis is its clinically
canonical instrument (CTQ, PSQI, ISF, BMI, CVLT, FAST, then the short substance and
mania scales). The same framework exposes where existing cohorts are currently
under-measured (Fig.~\ref{fig:voi}, right): schizophrenia participants were assessed on
markedly fewer immunometabolic items than bipolar participants ($20.7$ vs.\ $34.9$ items
on average) and fewer sleep items ($3.9$ vs.\ $8.4$), while depression is thinnest on
substance-use items---precisely the axes on which each cohort's under-measurement is
most consequential. For a consortium harmonising protocols across disorders, this
converts the measurement model into an actionable specification---a common core battery
plus disorder-specific top-ups---rather than a post-hoc analysis. The item ranking uses
pooled item parameters, so it assumes a common ordering across cohorts; the gaps
themselves are empirical. The arithmetic that quantifies a coordinate's uncertainty
therefore also prescribes how to reduce it most cheaply---a concrete clinical and
consortium-level use of the map beyond description.

\begin{figure}[tbp]
  \centering
  \includegraphics[width=\linewidth]{fig_value_of_information.png}
  \caption{\textbf{Value of information: what a future cohort should collect.}
  \textbf{Left}: greedily adding the single most-informative item across all eight axes
  (exact Fisher information per likelihood family), a shared $27$-item battery brings
  every axis online at a mean reliability of $0.70$; the first item chosen on each axis
  is its canonical instrument. \textbf{Right}: mean items collected per axis by cohort,
  shaded relative to the best-measured cohort. Schizophrenia is under-measured relative
  to bipolar disorder on immunometabolic ($20.7$ vs.\ $34.9$ items) and sleep ($3.9$ vs.\
  $8.4$) indicators, and depression is thinnest on substance-use items---the highest-value
  additions for future data collection at each site.}
  \label{fig:voi}
\end{figure}

\subsection{Designing efficient, harmonised batteries}

Because the model quantifies each indicator's information contribution---using the
correct form for each measurement family---and tolerates arbitrary missingness, it
functions as a design tool for future data collection, not only a description of the
current one. For a consortium harmonising protocols across bipolar disorder,
schizophrenia and depression, the shared battery and the cohort-gap map together
convert the measurement model into an actionable specification: a common core battery
(Fig.~\ref{fig:voi}, left) plus disorder-specific top-ups where a cohort is currently
blind (Fig.~\ref{fig:voi}, right), rather than a post-hoc analysis. The item ordering
uses pooled item parameters; a cohort-stratified refit, left to future work, would let
the most-informative set itself differ by diagnosis.

\subsection{The eight dimensions of the transdiagnostic atlas}\label{subsec:eightdim}

From ten candidate constructs, one global model returned a stable eight-dimension
structure: a general functional-burden factor (\Gfac) plus seven weakly correlated
specific axes---cognition, immunometabolic, sleep, mania, suicidality, developmental
risk and substance use (Figure~\ref{fig:atlas}; Table~\ref{tab:dimensions}).
Inter-factor correlations among the specific axes were weak throughout (mean absolute
off-diagonal $\approx 0.08$; Figure~\ref{fig:factors}b), confirming genuinely distinct
axes rather than one collapsed dimension; the single notable coupling was a modest
mania--sleep correlation ($\approx 0.24$). \Gfac{} was anchored only by
functioning and global-severity indicators (FAST loading $0.78$, EGF $0.60$) and
carried no symptom content, supporting its conservative reading as transdiagnostic
\emph{functional burden} rather than a symptom-liability ``\textit{p}'' factor.
Depression and anxiety instruments (MADRS, QIDS, STAI) did not form a separate factor;
they loaded $0.61$--$0.82$ on \Gfac{} as cross-loading ``windows'' onto burden. The
biology axis is a single \emph{immunometabolic} factor---cardiometabolic and
inflammatory markers (body-mass index, $\mathrm{HbA_{1c}}$, lipids, C-reactive protein)
load on one common dimension, anchored by body composition (BMI loading $\approx 0.95$,
CRP $\approx 0.37$; all $37$ loadings with credible intervals in
Table~\ref{tab:immet_loadings}), the data placing cardiometabolic and immune burden on a
shared axis rather than two. Body size anchors this axis but does not constitute it:
excluding the three anthropometric indicators and refitting, a coherent
cardiometabolic--inflammatory axis survived---all $34$ remaining markers credible, re-anchored
on blood pressure, lipids and inflammation, and highly congruent with the full axis (Tucker
$\varphi=0.90$; Supplement~\ref{annA:nobmi}). Of the ten
prior candidate factors, anhedonia was rejected (absorbed into \Gfac{} and the
depression/anxiety windows), and a further three constructs with no common
indicator---impulsivity, negative symptoms and sensory abnormality---were declared
\emph{not testable}, each verdict an outcome of the analysis rather than an assumption
(Table~\ref{tab:dimensions}). The structure was not an artefact of the priors: an independent flat-prior refit reproduced
the loadings and factor correlations essentially exactly (Tucker congruence $\varphi = 1.00$).
Information criteria also ranked the bifactor above unidimensional and correlated-factor
alternatives, but we treat this as secondary evidence, since such criteria are known to favour
bifactor structures for reasons of model form independent of the data-generating
structure~\citep{bonifaycai2017,greene2019fit}. Absolute fit was
acceptable across both likelihood blocks (continuous Bayesian standardized
root-mean-square residual $\approx 0.07$; $21/22$ non-Gaussian indicators reproduced
within the $90\%$ posterior-predictive interval). \emph{Clinically,} the axes answer two
different questions: \Gfac{} answers \emph{how impaired} the patient is, while the seven
specific axes answer \emph{what kind} of burden the patient carries.

The map is otherwise simple-structure---each indicator loads on its home axis and
\Gfac{}---by design rather than by fiat. A sparsity-inducing regularized (``Finnish'')
horseshoe prior is placed on every off-home specific-to-specific loading: a global
shrinkage term pulls the whole set toward zero (the prior's default belief is ``no
cross-loadings''), heavy local tails let a genuinely supported cross-loading escape,
and a slab caps its magnitude. This protects the thin factors (substance, mania, each
defined by very few indicators) from dilution while letting small, clinically real
cross-talk emerge. Given total freedom---a continuous sparse-ESEM refit that frees
every off-home cell---$\approx 83\%$ shrank to $\approx 0$ (the simple structure is
\emph{recovered from the data}, not imposed), and exactly three cross-loadings survived into the full
model with $95\%$ credible intervals excluding zero, all onto cognition: childhood
adversity (CTQ, $-0.094$) and two sleep items (PSQI latency $+0.057$, daytime
dysfunction $-0.070$). Given the freedom to load anywhere, the model reproduced known
clinical cross-talk---adversity and poor sleep weighing weakly on cognition---and
nothing spurious: positive evidence that the map is well-specified, not an artefact of
rigid zeros (Supplement~\ref{ann:measurement}).

\begin{figure}[!tp]
  \centering
  \includegraphics[width=\linewidth]{fig2_map.png}
  \caption{\textbf{Posterior loading atlas of the measurement model.} Rows are clinical
  indicators grouped by their home factor (coloured left strip and alternating bands);
  columns are the general functional-burden dimension \Gfac{} and the seven specific
  latent axes (D1--D7). Dot size and colour encode the absolute posterior median loading;
  outlined dots indicate loadings whose $95\%$ credible interval excludes zero, and a
  heavier ring marks the home-factor anchor. Specific axes are each anchored by a
  coherent clinical indicator block, whereas \Gfac{} captures broad functional burden
  and carries cross-loading ``windows'' over the depression and anxiety instruments
  (MADRS, QIDS, STAI; shaded on the \Gfac{} column). The three data-supported off-home
  cross-loadings---childhood adversity and two sleep items onto cognition---render as
  outlined dots in the cognition column. The full eight-factor map is shown
  (continuous backbone plus the mania and substance axes from the joint mixed fit); the
  complete atlas over all modelled indicators is \edfig{fig:fullatlas}. This structure
  supports an atlas-like reading of the latent space: specific dimensions are clinically
  localised axes, while \Gfac{} acts as a transdiagnostic burden backbone. Factor-wise
  interpretability and the inter-factor correlation matrix are in
  Figure~\ref{fig:factors}.}
  \label{fig:atlas}
\end{figure}

\begin{figure}[tbp]
  \centering
  \includegraphics[width=\linewidth]{fig_factors.png}
  \caption{\textbf{Factor interpretability and inter-factor correlations.}
  (\textbf{a}) Factor-wise lollipops: for each axis (\Gfac{} and D1--D7), its strongest
  home indicators with $95\%$ credible whiskers---``what defines each axis.''
  (\textbf{b}) The inter-factor correlation matrix $\Phi$ is weak throughout
  (specific--specific mean $|\text{off-diagonal}|\approx 0.08$); \Gfac{} is orthogonal to
  the specific axes by construction (top row and column). The one notable specific
  coupling is mania--sleep ($\approx 0.24$); substance is pinned orthogonal to the
  correlated block (its cross-factor correlations are non-identifiable), and the
  biology$\times$symptom cells all stay at the off-diagonal floor.}
  \label{fig:factors}
\end{figure}

\subsection{How each domain relates to general burden}\label{subsec:biologyg}

Because the bifactor identification fixes \Gfac{} orthogonal to the specifics by
construction, a domain's relationship to general burden is read from its indicators'
\emph{direct loadings on \Gfac{}} (Figure~\ref{fig:factors}). On this metric three axes
sit at the periphery, loading only weakly on \Gfac{}: immunometabolic ($\approx 0.06$),
substance ($\approx 0.07$) and developmental-risk ($\approx 0.07$)---a near-tie that, by
itself, singles out none of them. To convert ``biology is separable from severity'' from a
modelling constraint into a measured quantity, we re-estimated the model allowing \Gfac{}
to correlate freely with the specific axes. The resulting vector of general--specific
correlations, $\phi_{Gd}=\operatorname{Corr}(\Gfac,D)$ (Methods, Eq.~\ref{eq:corrg}), is
the \emph{primary estimand} of FACE-ATLAS. Under this correlated-\Gfac{} model,
immunometabolic load was the \emph{least} entangled with functional burden: a correlation
with \Gfac{} of $\approx 0.10$, against $0.39$ for cognition and $0.42$ for sleep
(Figure~\ref{fig:biologyg}). In shared-variance terms the gap is stark---immunometabolic
load shares only $\approx 1\%$ of its variance with general burden, against $\approx 15\%$
(cognition) and $\approx 18\%$ (sleep).

What distinguishes immunometabolic load from the other two peripheral axes is that it is
decoupled from the \emph{entire} clinical picture, on a freely-estimated basis. Substance
is orthogonal \emph{by construction} (pinned, because its cross-factor correlations are
non-identifiable), so its independence is imposed rather than measured; developmental-risk,
though weak on \Gfac{}, couples to the symptom axes (sleep and suicidality,
$\Phi\approx 0.20$) and indexes a historical antecedent rather than a current state.
Immunometabolic load alone is weakly coupled both to functional burden \emph{and} to the
symptom axes---its largest inter-factor correlation is $\approx 0.08$, against $\ge 0.16$
for every other freely-correlating factor (Figure~\ref{fig:factors})---and is, uniquely, an
earned ``island'' in the map. (We note one boundary: the correlated-\Gfac{} arm cleanly
separates biology from cognition and sleep but not from developmental-risk, so the
immunometabolic--developmental distinction rests on the symptom-decoupling in $\Phi$, not on
$\phi_{Gd}$; Supplement~\ref{ann:measurement}.) Two patients with the same overall
impairment can therefore carry very different immunometabolic loads---a patient's biological
burden rides largely on an axis that the global clinical-severity picture does not capture.
This dissociation was robust to the obvious
confounders. Adjusting each indicator for age, sex, education and recruitment site, and
then additionally for antipsychotic exposure, did \emph{not} raise the immunometabolic
correlation---if anything it lowered it (to $\approx 0.07$), and the immunometabolic axis
remained the least entangled at every step (Figure~\ref{fig:biologyg};
Supplement~\ref{ann:measurement}). Consistent with a large independent literature
(Discussion), we state this as relative, not absolute, independence: immunometabolic load
is \emph{largely independent of a general functional-burden axis}, the low correlation
quantifying---at the low end of what that literature would predict---a pattern previously
described only qualitatively. \emph{Clinically,} this means two patients with comparable
functional burden can have sharply different cardiometabolic and inflammatory profiles:
severity cannot be used as a proxy for biological load. This is the property a
biology-aware stratification can exploit, and it is one the downstream stages
carry forward and re-test rather than assume.

\begin{figure}[t]
  \centering
  \includegraphics[width=0.80\linewidth]{fig3_biology_g.png}
  \caption{\textbf{Immunometabolic load is the least burden-entangled domain.}
  Correlation of each continuous specific axis with the general
  functional-burden factor \Gfac{} (bars), unadjusted and after adjustment for
  medication, adiposity and site; the bifactor general loading $|\lambda_G|$
  (diamonds) is shown for comparison. Immunometabolic load falls in the
  ``biology zone'' ($\Phi<0.15$) and stays there under adjustment, whereas cognition
  and sleep partly track \Gfac{}.}
  \label{fig:biologyg}
\end{figure}

\subsection{The coordinate cloud is continuous, not discretely clustered}

We next asked whether these eight coordinates fall into separable kinds. They do not---a
\emph{negative structural result}. On the cloud $\{(\hat f_i,S_i)\}$, a
structure-discovery battery run over the posterior
draws---cluster tendency, modality, silhouette, the gap statistic, a Gaussian-mixture
information criterion, density-based clustering and a topological summary---returned a
continuum verdict: the silhouette peaked at only $0.146$ (below the $0.15$ separation
floor), the gap-statistic optimum was a single component, density-based clustering
assigned $100\%$ of points to noise, and the topological summary
was a single connected component. A single-Gaussian \emph{falsification null} made the
verdict decisive: the best partition of the real cloud separated patients no better than
slicing a structureless Gaussian blob of identical covariance (best silhouette $0.140$
for the real data versus $0.137\pm0.002$ for the null, $z=1.13$, n.s.; the
Gaussian-mixture-optimal number of components collapsed to one in all $20$ posterior
draws; Figure~\ref{fig:continuum}a). There are, in this clinical--biological space, no
well-separated, reproducible discrete clusters. This is a statement about discrete clustering
in this representation (Gaussian-copula ESEM coordinates at $N=9{,}013$ under this pipeline),
not a claim that biological subtypes are absent in other feature spaces or larger
samples~\citep{marquand2016,parker2025,haslam2012}, nor a denial of the clinical heterogeneity
we do find.

The continuum is genuine and transdiagnostic, not an artefact. Diagnosis is intermixed
throughout the cloud (the soft tessellation's adjusted Rand index is $0.006$ against the seven
DSM-5 subtypes and $0.011$ against cohort, both $\approx 0$; the $A=5$ archetype assignment is
likewise near-orthogonal, ARI $0.021$, \S\ref{subsec:archetypes}; Figure~\ref{fig:continuum}b); overall burden
forms a smooth gradient along the principal axis (Figure~\ref{fig:continuum}c); and
immunometabolic load varies along a \emph{different} direction
(Figure~\ref{fig:continuum}d)---the biology--\Gfac{} dissociation, made visible.
The verdict is not a missingness artefact (a classifier predicting position from each
patient's coverage pattern performed below baseline, accuracy $0.611$ versus a $0.683$
baseline, lift $-0.072$) and is a tighter description of the coordinate cloud than
diagnosis (free-mixture information criterion $185{,}600$ versus $188{,}200$ for a
DSM-5-constrained model, $\approx 3\times$ the explained variance at lower BIC). These
same properties explain why clustering
the \emph{raw} indicators can appear to yield biotypes where the latent space does
not~\citep{pan2026}: a raw partition is driven largely by measurement-side
structure---structured missingness, a few heavy-tailed laboratory variables, redundant
indicator blocks, and site-specific instruments that re-encode diagnosis---which the
measurement model neutralizes by construction (observed-cell likelihood, per-indicator and
Gaussian-copula likelihoods, shared factors, and diagnosis held out under invariance), so
the structure that remains is genuinely transdiagnostic and tested against a structureless
null rather than assumed. \emph{Clinically,} the appropriate object is
therefore a coordinate with uncertainty (Eq.~\ref{eq:res-post}), not a categorical
biotype.

\begin{figure}[t]
  \centering
  \includegraphics[width=0.99\linewidth]{fig4_continuum.png}
  \caption{\textbf{The patient space is a continuum, not biotypes.} (\textbf{a}) A
  single-Gaussian falsification null: the best silhouette achievable on the real
  eight-dimensional cloud is no better than on a structureless Gaussian of matched
  covariance ($z=1.13$, n.s.); the mixture-optimal number of clusters is one in all
  posterior draws. (\textbf{b--d}) Principal-component embedding of the coordinates
  ($N=9{,}013$), coloured by diagnosis (fully intermixed), by general burden \Gfac{}
  (a smooth gradient), and by immunometabolic load (a gradient in a \emph{different}
  direction).}
  \label{fig:continuum}
\end{figure}

\subsection{The measurement map is stable to holding out recruitment sites}

The map is estimated once on the full sample; a reviewer may reasonably ask
whether its structure is an artifact of the particular sites that contributed
patients. Because the cohort spans 21 expert centres, we can test this directly:
for each of the 15 sites contributing at least 100 patients, we refit the entire
measurement model with that site's patients removed (held-out $N$ ranging from
134 to 1{,}308) and compared the leave-site-out map to the full-data map by
Tucker's congruence coefficient per factor. Refitting used the variational
path, which reproduces the reference posterior loadings at Tucker
$0.957$--$0.999$ (Annex~H), making 15 refits computationally tractable.

The immunometabolic axis---the load-bearing structure of this paper---was
essentially invariant to which site was removed: loading congruence with the
full-data map was $\varphi \geq 0.999$ in every fold (minimum $0.9993$, median
$0.9998$). Its dissociation from general severity, measured at the score level
as $|\mathrm{corr}(\text{immunometabolic},\ \text{general burden})|$, remained
between $0.073$ and $0.082$ across all folds (full-sample score-level value $\approx 0.08$),
and the immunometabolic axis was the least severity-correlated of the seven
specific factors in all 15 of 15 refits. The two single-diagnosis holdout sites
(one bipolar-only, one schizophrenia-only) behaved no differently from the
multi-diagnosis sites. Across all eight dimensions, every factor cleared its
pre-registered congruence bar (Tucker $\geq 0.95$ for the four major axes,
$\geq 0.85$ for the four thin specific factors) in all 15 folds; the single
lowest per-factor value across the entire analysis was $\varphi = 0.917$, for a
two-to-four-item thin factor at a small held-out site ($n=237$, Monaco), still comfortably
above its bar (\edfig{fig:loso}).

\subsection{Five archetypal extremes define the shape of the continuum}\label{subsec:archetypes}

A continuum has no clusters, but it still has a \emph{shape}. Archetypal analysis recovers
that shape as a small set of reproducible extreme corners---a \emph{positive geometric
result} that coexists with the negative clustering result above. Each patient is written
as a convex blend of $A=5$ extreme phenotypes,
\begin{equation}
  \hat f_i\ \approx\ \sum_{a=0}^{4} w_{ia}\,z_a,\qquad
  w_{ia}\ge 0,\qquad \sum_{a=0}^{4} w_{ia}=1 .
  \label{eq:res-archetype}
\end{equation}
The five $z_a$ are extreme points on the convex hull of the patient cloud, and the
weights $w_{ia}$ are \emph{barycentric coordinates}: they quantify how strongly each
patient is pulled toward each extreme. Dominant-archetype labels are used only for
summaries; the inferential object is the continuous weight vector $w_i$, not a class
assignment. The number of corners ($A=5$) was set by cross-seed stability rather than by
a parsimony rule---five corners were the largest reproducible solution (cross-seed Tucker
congruence $0.979$, explained variance $0.60$), with a clean stability cliff at six
corners ($0.979\!\to\!0.436$)---and uncertainty was propagated by refitting
Eq.~\eqref{eq:res-archetype} across posterior draws (Supplement~\ref{ann:stratification}).

The five extremes separate biology from symptoms from severity (Figure~\ref{fig:archetypes};
Table~\ref{tab:archetypes}):
an \emph{activation / sleep} corner (A0; high sleep and mania), a \emph{severe,
clean-biology} corner (A1; high overall severity but low immunometabolic load), an
\emph{immunometabolic} corner (A2; high immunometabolic load with severity and
suicidality)---the biology pole---a \emph{trauma / suicidality} corner (A3; high
developmental risk and suicidality), and a \emph{low-burden} pole (A4; everything low).
The immunometabolic corner A2 and the severe corner A1 carry comparable general burden
but opposite biological profiles---A1 versus A2 is precisely why \emph{severity is not
biology}, the central dissociation now rendered as geometry. The symptom space itself
resolves into two distinct corners---activation/sleep (A0) and trauma/suicidality
(A3)---rather than one combined symptom pole. Patients sit \emph{between} the corners far
more than at them: the weight vectors are diffuse (median normalized entropy $0.65$), so
most patients are convex blends rather than single-corner cases. Critically, because
immunometabolic load varies continuously without a density \emph{gap}, it never forms a
clean partition boundary at any granularity; it surfaces instead as a \emph{direction of
maximal phenotype}---an archetypal corner---which is exactly why the continuous
coordinates and the weight vectors $w_i$, not any hard tessellation, are the primary
objects downstream. \emph{Clinically,} the archetypes provide a vocabulary for reading the
continuum, not a replacement for it.

Read against the held-out diagnostic labels, the archetype assignment is essentially
orthogonal to DSM-5 (adjusted Rand index $0.021$) and every corner is populated in all three
cohorts---the stratification analogue of diagnosis-as-metadata. Each diagnosis is a
\emph{mixture} of corners and each corner draws from every diagnosis (Table~\ref{tab:archcomp}):
the activation/sleep corner is bipolar-enriched (mania is a bipolar feature) and the biology
corners (A1/A2) are depression/schizophrenia-enriched, yet none is diagnosis-specific. The
simplex nonetheless describes \emph{where a patient sits} on the eight-dimensional map far more
tightly than the diagnostic label: averaged over the eight axes, the archetype assignment
accounts for $\eta^2=0.256$ of the coordinate variance against only $0.026$ for the DSM-5
grouping---an $\approx\!9.7$-fold improvement, at lower BIC. A categorical diagnosis is thus a
weak organiser of a patient's clinical--biological state and the archetype mixture an informative
one, transdiagnostically; whether the mixture also \emph{predicts} better, beyond being a tighter
\emph{description}, is the prognostic question taken up below.

Stated precisely, the archetypes separate on severity \emph{and} profile. Severity (\Gfac) is in
fact the single strongest separating axis ($\eta^2=0.47$)---the simplex spans a clear gradient
from the well pole (A4) to the severe poles (A1/A2)---but the corners separate almost as strongly
on immunometabolic load ($0.39$), sleep ($0.38$) and developmental risk ($0.35$). This is the
value over a one-number severity score: at any given severity the archetype distinguishes patients
by biological and symptom \emph{profile}, so A1 and A2---equally severe, biologically opposite---are
different corners a severity score cannot tell apart. (The sharper ``driven by activation and
suicidality more than severity'' statement belongs to the coarse two-region tessellation,
$\eta^2\approx0.22$ each versus $0.05$ for \Gfac{}; a distinct, lower-resolution object,
Supplement~\ref{ann:stratification}.) Choosing soft archetypal corners over hard clustering is
itself the methodological point: on a cloud with no density gaps, convex mixtures of extreme
phenotypes neither manufacture boundaries the data lack nor blur patients into centroid
``averages,'' and they preserve a patient's position relative to \emph{every} pole rather than
collapsing it to a bin.

\begin{table}[t]
\centering
\caption{\textbf{Archetype composition by cohort (diagnosis-blind).} Share of each cohort's
patients whose dominant ($\arg\max$) archetype is each corner. Every corner is populated in all
three cohorts and each diagnosis is a mixture of corners; the archetype assignment is essentially
orthogonal to DSM-5 (adjusted Rand index $0.021$).}
\label{tab:archcomp}
\small
\renewcommand{\arraystretch}{1.15}
\begin{tabular}{@{}l ccc@{}}
\toprule
Dominant corner & BP & DR & SZ \\
\midrule
A4 low-burden / well     & 0.30 & 0.11 & 0.29 \\
A3 trauma / suicidality  & 0.24 & 0.22 & 0.18 \\
A0 activation / sleep    & 0.19 & 0.06 & 0.10 \\
A1 severe, clean-biology & 0.14 & 0.31 & 0.25 \\
A2 immunometabolic       & 0.13 & 0.30 & 0.19 \\
\bottomrule
\end{tabular}
\end{table}

\begin{figure}[t]
  \centering
  \archetypesimplexblend
  \caption{\textbf{The archetype simplex: every patient a convex blend of five extremes.}
  The five archetypal extremes are the corners---\textbf{A0} activation / sleep,
  \textbf{A1} severe but clean-biology, \textbf{A2} immunometabolic (the biology pole),
  \textbf{A3} trauma / suicidality, and \textbf{A4} low-burden---and every patient sits
  \emph{inside} the simplex as a convex blend whose barycentric weights $w_i$ sum to one
  (worked example, foot). The central dissociation is geometric: \textbf{A1} and
  \textbf{A2} are distinct corners at essentially the \emph{same} general burden but
  opposite biology---the biology--severity dissociation, rendered as a direction one can point to.}
  \label{fig:archetypes}
\end{figure}

\paragraph{The archetype simplex is a fog, not five clusters.} The continuum verdict is equally
visible in archetype space. Writing each patient as their membership vector $w_i\in\Delta^{4}$
over the five corners and measuring how \emph{blended} they are by the entropy
$H(w_i)=-\sum_{a}w_{ia}\log w_{ia}$, the $9{,}013$ patients fill the interior of the simplex
rather than piling at its vertices (Fig.~\ref{fig:simplexfog}): the median entropy is $1.24$ of a
possible $\log 5 = 1.61$ nats, $74\%$ of patients have no majority archetype (largest weight
$<0.5$), and only $0.4\%$ sit near a single corner (largest weight $>0.8$). The archetypes are
extremes that \emph{span} a continuous cloud, not labels that \emph{partition} it---the same
conclusion the structure tests reach, now read directly off the membership weights.

\begin{figure}[tbp]
  \centering
  \includegraphics[width=\linewidth]{fig_simplexfog.png}
  \caption{\textbf{The archetype simplex is a fog, not five clusters.} \textbf{(A)}~Each of the
  $9{,}013$ patients placed at the barycentre of the five archetype corners by its membership
  weights $w_i$, coloured by the entropy $H(w_i)$ of those weights; the cloud fills the interior
  rather than the vertices. \textbf{(B)}~Distribution of $H(w_i)$: median $1.24$ of a maximum
  $\log 5 = 1.61$ nats, with $74\%$ of patients having no majority archetype and only $0.4\%$ near
  a single corner---the archetypes span the continuum, they do not partition it.}
  \label{fig:simplexfog}
\end{figure}

\subsection{The geometry persists over two years (the map persists)}

Scored forward onto the $12$- and $24$-month visits on the \emph{fixed} model
(observed cells, uncertainty propagated, never re-estimated; baseline coordinates
reproduced at Pearson $r\approx0.99$ through this path), the measurement remained
invariant for all four testable backbone axes (longitudinal Tucker congruence
$\varphi = 0.987$--$0.996$; the immunometabolic axis invariant at $\varphi=0.987$;
\edfig{fig:invariance}). Coordinate change can therefore be read as patient change rather
than instrument drift. A measurement-error decomposition then separated trait from state
(Figure~\ref{fig:traitstate}). The durable, trait-like axes were the biological and
cognitive ones, with immunometabolic load the single most durable axis (test--retest
intraclass correlation: immunometabolic $0.91$, cognition $0.70$); the moving, state-like
axes were the symptom ones (substance $0.49$, sleep $0.47$, suicidality $0.43$,
developmental risk $0.39$). General burden was trait \emph{by rank} (intraclass
correlation $0.62$): patients largely kept their relative position even as the cohort as a
whole improved on severity and symptoms (population shift $-0.46$ on \Gfac{} and $-0.84$
on suicidality, while immunometabolic load stayed static, $+0.04$). The archetype geometry
persisted---each patient's weight vector $w_i$ moved little over two years (median cosine
similarity $0.81$). \emph{Clinically,} the reading is compact: \emph{stratify on the
durable biology; monitor the moving symptoms}.

\begin{figure}[t]
  \centering
  \includegraphics[width=0.98\linewidth]{fig5_persistence.png}
  \caption{\textbf{Durable biology, moving symptoms.} (\textbf{a}) Trait fraction
  (test--retest intraclass correlation, with $94\%$ credible intervals) over
  follow-up: the biological and cognitive axes are trait-like, the symptom axes
  state-like; mania is data-limited (hatched). (\textbf{b}) Individual rank stability
  versus population-level change: patients keep their biological rank (high intraclass
  correlation, near-zero population shift) while the cohort slides down the severity
  and symptom axes.}
  \label{fig:traitstate}
\end{figure}

\paragraph{A trait/state spectrum across the axes.} The same decomposition, shown in full as a
variance \emph{thermometer} (Fig.~\ref{fig:thermometer}), partitions each axis into trait, state
and measurement error and so sharpens the reading. Ordered by error-corrected trait fraction the
eight axes form a single durability spectrum from immunometabolic biology (ICC $0.91$; $80\%$ of
its total variance is the stable between-patient trait) down to developmental risk ($0.39$), with
cognition ($0.70$) and general burden ($0.62$) trait-dominant and the symptom axes mixed.
Separating state from measurement error also shows that the symptom axes (sleep, suicidality,
developmental risk) carry genuine within-patient change rather than mere noise, whereas mania's
variance is dominated by measurement error---its two indicators leave it data-limited and its
trait verdict is not relied upon. The biology that anchors the map is again its most durable
dimension, exactly the property a stratification axis should have.

\begin{figure}[tbp]
  \centering
  \includegraphics[width=0.92\linewidth]{fig_traitstate.png}
  \caption{\textbf{Trait/state thermometer.} Each axis's variance across the three visits,
  partitioned into a stable between-patient (trait) part, a moving within-patient (state) part and
  measurement error, normalised to the axis total and ordered by the error-corrected trait
  fraction (ICC $=$ between/(between$+$within)). The immunometabolic biology axis is the most
  durable (ICC $0.91$), developmental risk the most state-like ($0.39$); mania is mostly
  measurement error (two indicators, data-limited). This variance view complements
  Fig.~\ref{fig:traitstate}, separating genuine within-patient state from measurement noise.}
  \label{fig:thermometer}
\end{figure}

\subsection{Archetype position predicts group-level two-year functioning (the map predicts)}

The clinically actionable prognostic signal sits at the level of \emph{position on the
map}, not individual point prediction, and it is visible before any model-comparison
statistic. Ranking patients by their dominant archetypal extreme, two-year functional
remission forms a gradient that reproduces \emph{within} every diagnosis: the
immunometabolic corner (A2) is the worst- or near-worst-remitting archetype in bipolar
disorder, schizophrenia and depression alike (bipolar $27\%$ vs.\ $73\%$ at the
low-burden pole A4; schizophrenia $9\%$ vs.\ $25\%$; depression $31\%$ vs.\ $72\%$;
Figure~\ref{fig:prognosis}a)---so the gradient is a property of the map, not a
cohort-composition artefact. Pooled across cohorts the same ranking spans
$22\%$--$63\%$ (Figure~\ref{fig:glance}), a secondary, marginal framing reported only
alongside the within-cohort result. The archetype coordinate is also a far more compact
descriptor of where a patient sits on the map than diagnosis: it accounts for a mean
$\eta^2$ of $0.256$ across the eight axes, against $0.026$ for DSM-5 category and
$0.018$ for cohort---an $\approx9.7$-fold tighter summary (Figure~\ref{fig:prognosis}b).
This is compactness on the map's own axes, not a claim about outcome variance
explained; the outcome evidence is the remission gradient above, kept as a separate
statement.

We then asked whether a baseline coordinate or phenotype predicts this two-year
functioning outcome \emph{incrementally beyond} diagnosis, current severity and the
baseline value of the outcome, using an errors-in-variables Bayesian model that
propagated each coordinate's measurement uncertainty. The
map added predictive signal for \emph{functioning} but not for severity (which was
saturated by its own baseline autoregression). The signal was carried by the
\emph{phenotype}, not by any single laboratory axis: the five archetypal extremes
added the most held-out predictive value (expected log predictive density gain $+62.8$
over the diagnosis-plus-severity-plus-baseline reference; Table~\ref{tab:incremental}),
the seven specific axes added $+38.1$, while the isolated ``durable biology'' axes
\emph{alone} sat at the bar ($+2.3$, ambiguous). Biology still predicts functioning,
but as a \emph{corner of the phenotype} (the immunometabolic archetype) rather than as
a standalone block---the immunometabolic axis itself carries a credible adverse
direction (higher load $\to$ worse future functioning). No hard partition improved on
the continuous and archetype encodings (every tessellation added $\approx+16$ to
$+20$, strictly dominated by the archetypes); the operative number of clusters is
therefore \emph{none}---the actionable object is a continuous position, not a category.
The map was \emph{co-informative} with DSM-5 rather than redundant with it
(functioning: map $+17.3$, DSM-5 $+29.0$, both $+62.6$; Figure~\ref{fig:prognosis}c),
and \emph{course-dependent}, with the signal concentrated in the episodic, open-course
cohorts (bipolar disorder and depression) and weak when bipolar disorder was removed.
At the \emph{individual} level, the same map adds only $+0.010$ to remission-prediction
AUC over a diagnosis-plus-severity reference---the map's prognostic value is one of
stratification and durable measurement (test--retest ICC $0.91$), not a replacement for
individual clinical judgment.

\begin{table}[t]
\centering
\caption{\textbf{Incremental predictive value for the two-year outcome ($\Delta$ELPD).}
Held-out expected log predictive density gain of each baseline encoding when added to the
reference model (diagnosis $+$ baseline severity $+$ baseline outcome autoregression), for
functioning (EGF) and clinician severity (CGI-S). The archetype encoding is the strongest
predictor of \emph{functioning}; the same encoding gains little on \emph{severity} (saturated
by its own baseline); the continuous and archetype encodings dominate every hard tessellation
(operative number of clusters $=$ none).}
\label{tab:incremental}
\small
\renewcommand{\arraystretch}{1.15}
\begin{tabular}{@{}l c c@{}}
\toprule
Encoding added to the reference & $\Delta$ELPD functioning & $\Delta$ELPD severity \\
\midrule
$+$ $A{=}5$ archetypes (all axes)        & $\mathbf{+62.8}$ & $+11.9$ \\
$+$ continuous coordinates (8 axes)      & $+38.1$ & $+13.3$ \\
$+$ $A{=}5$ archetypes ($\perp\Gfac$)    & $+33.5$ & $+14.0$ \\
$+$ tessellation $K=3\,/\,4\,/\,2$       & $+19.6\,/\,16.6\,/\,15.9$ & $+5.6\,/\,4.4\,/\,4.9$ \\
$+$ durable (immunometabolic) axis       & $+2.3$ (amb.) & $-0.7$ (amb.) \\
reference (autoregression)               & $0$ & $0$ \\
\bottomrule
\end{tabular}
\end{table}

This within-cohort gradient (Figure~\ref{fig:prognosis}a) is a genuine within-diagnosis
effect, not a cohort-composition artefact: the rank held inside every cohort
(immunometabolic corner $\to$ well pole: bipolar $27\%\!\to\!73\%$, depression
$31\%\!\to\!72\%$, schizophrenia $9\%\!\to\!25\%$; a cohort-adjusted best-versus-worst
odds ratio $\approx 6.3$), standardizing to a common cohort mix moved it by only
$\approx 4\%$, and the corner-by-cohort interaction was non-significant. A2 also
remained worse than the equally severe clean-biology corner A1 within diagnosis ($22\%$
vs $34\%$ pooled)---the severity-is-not-biology dissociation, now translated into
outcome. A wide spread in \emph{group} remission rates alongside the small
\emph{individual} AUC gain noted above is the expected signature of groups that differ
in average risk but overlap substantially at the patient level, layered on an
already-strong autoregressive baseline---the map enriches and forecasts at the group
level without sharply separating individuals. The functioning signal survived
inverse-probability weighting for attrition ($+54.4$),
restriction to either two of the three cohorts ($+56.4$, $+61.8$), and a permutation null
(which correctly vanished). Finally, against the full set of raw indicators under a
matched gradient-boosting classifier, the eight-dimension map was a \emph{sufficient}
representation for predicting deterioration (equal AUC) and \emph{near-sufficient} for
recovery (raw indicators $+0.04$ AUC), with $92$--$97\%$ of the raw recovery signal
already contained within the eight factors (\edfig{fig:repbench})---the compression is
faithful, not lossy in any axis the map omits. \emph{Clinically,} the atlas is useful for
group-level enrichment and monitoring, not as an individual remission calculator.

\begin{figure}[t]
  \centering
  \includegraphics[width=0.99\linewidth]{fig6_prognosis_rebuilt.png}
  \caption{\textbf{The archetype coordinate carries outcome structure diagnosis
  misses.} (\textbf{a}) The within-cohort remission gradient: two-year functional-remission
  rate by baseline dominant archetype, shown separately within each cohort. The
  immunometabolic pole (A2) is the lowest-remitting archetype in every diagnosis (bipolar
  $27\%$ vs.\ $73\%$ at the low-burden pole A4; schizophrenia $9\%$ vs.\ $25\%$; depression
  $31\%$ vs.\ $72\%$)---a within-diagnosis gradient, not a cohort-composition artefact.
  (\textbf{b}) The archetype assignment summarises the eight map axes $9.7\times$ more
  compactly than DSM-5 (mean $\eta^2$ across the eight axes: archetype $0.256$, DSM-5
  diagnosis $0.026$, cohort $0.018$); this is compactness on the map's own axes, not
  outcome variance explained. (\textbf{c}) Held-out incremental predictive value
  ($\Delta$ELPD) for two-year functioning: the archetype encoding is the strongest single
  predictor, ahead of the continuous coordinates and every hard tessellation. The
  individual-level remission AUC gain over diagnosis is small ($+0.010$); the map's
  prognostic value is in stratification and durable measurement, not individual point
  prediction.}
  \label{fig:prognosis}
\end{figure}

\medskip\noindent\textbf{Treatment moderation: a bounded, defensive result.}
On observational treatment-as-usual---not randomised arms---we asked whether the map
\emph{moderates} treatment response, under an explicit bounds-and-defends design (overlap gate
$\to$ propensity adjustment $\to$ doubly-robust errors-in-variables moderation $\to$ E-value; full
analysis in Supplement~\ref{ann:treatment}). It does not reliably (Table~\ref{tab:moderation}):
the best-identified contrast, lithium in bipolar disorder, is a well-identified null (E-value
$1.20$--$1.28$, so a trivial unmeasured confounder would erase it); antipsychotic in bipolar
disorder shows a non-zero average effect but an E-value of only $1.80$ and a suggestive,
unconfirmed immunometabolic interaction; clozapine in schizophrenia is underpowered. With a
minimum detectable interaction of $\approx0.20$\,SD, the map is \emph{prognostic and descriptive,
not prescriptive} on these data. The analysis is nonetheless informative in one direction: it
\emph{defends} the prognostic result---adding the drug classes patients actually received to the
prognostic model attenuates the immunometabolic carrier by under $8\%$ (durable axis
$6.4\%/4.1\%$ under inverse-probability weighting; A2 corner $7.7\%/6.4\%$), so the forecast is
not a treatment proxy---and it \emph{describes} course: the immunometabolic corner faces the
hardest two-year course (treatment resistance $44\%$ and side-effects $25\%$ versus $20\%$ and
$11\%$ at the low-burden pole). Genuine treatment \emph{selection} would require randomised or
trial-arm data.

\begin{table}[t]
\centering
\caption{\textbf{Treatment moderation on observational treatment-as-usual (bounds-and-defends).}
For each within-cohort contrast: overlap (max standardised mean difference, before $\to$ after
weighting), average treatment effect on functioning (ATE), E-value (unmeasured-confounding
strength that would explain the effect away), the moderation $\Delta$ELPD (whether a
treatment $\times$ map interaction improves held-out fit), and the verdict. Full analysis in
Supplement~\ref{ann:treatment}.}
\label{tab:moderation}
\small
\renewcommand{\arraystretch}{1.15}
\begin{tabular}{@{}l c c c c >{\raggedright\arraybackslash}p{2.7cm}@{}}
\toprule
Contrast & overlap (SMD $\to$) & ATE & E-value & mod.\ $\Delta$ELPD & Verdict \\
\midrule
Lithium (BP)       & $0.30\!\to\!0.01$ & $-0.03$ (incl.\ 0) & $1.20$--$1.28$ & $-1.6$ & well-identified null \\
Antipsychotic (BP) & $0.71\!\to\!0.08$ & $-0.24$ (excl.\ 0) & $1.80$         & $+3.4$ & confounded; suggestive-only \\
Clozapine (SZ)     & $0.33\!\to\!0.07$ & $+0.03$ (incl.\ 0) & $1.21$         & $-1.8$ & underpowered ($n\!\approx\!515$) \\
\bottomrule
\end{tabular}
\end{table}

\subsection{A variational re-estimation reproduces the atlas at a fraction of the cost}

Because the audit-grade NUTS fit is expensive, we asked whether the same measurement model can
be recovered by a much cheaper estimator. Re-estimating the eight-factor map as a generalized
linear latent-variable model fit by stochastic variational inference (Methods;
Supplement~\ref{ann:variational}) reproduces the NUTS atlas on its two load-bearing objects at
roughly $100\times$ less cost (${\approx}4.5$ minutes on a CPU versus hours). The loading map is
the same map: loading columns are congruent on all eight factors (Tucker $\varphi$ $0.96$--$0.999$;
\Gfac{} $0.957$) and the full item$\times$factor loading scatter is $r=0.993$ with no sign flips
(Figure~\ref{fig:vinuts}a). Patients land in the same place: per-patient coordinates correlate
with the NUTS coordinates at $r\ge0.90$ on all eight axes ($\ge0.99$ for the
\Gfac{} and immunometabolic axes; Figure~\ref{fig:vinuts}b). The map is also a faithful generative
model---$9{,}013$ synthetic patients drawn from it match the observed marginals (Kolmogorov--Smirnov
median $0.04$) and covariance (correlation SRMR $0.078$, as for NUTS). The one quantity the
variational engine does \emph{not} reproduce is the inter-factor correlation $\Phi$: its structure
is correct (the biology$\,\perp\,$\Gfac{} zeros are exact) but the off-diagonals are systematically
attenuated by ${\sim}21\%$ under the mean-field posterior---a structural property of variational
inference, not a tunable defect, which is why NUTS is retained for correlation magnitudes and all
credible intervals (Supplement~\ref{ann:variational}). The variational engine is therefore a fast
operational and generative stand-in for the map and coordinates; a fuller development is left to
future work.

\begin{figure}[t]
  \centering
  \includegraphics[width=0.95\linewidth]{gllvm_oop/vi_vs_nuts_metrics.png}
  \caption{\textbf{A variational re-estimation reproduces the NUTS atlas at ${\sim}100\times$ less
  cost.} (\textbf{a}) Per-factor loading congruence (Tucker $\varphi$) between the variational
  GLLVM and the NUTS copula fit; all eight axes clear the major ($0.95$) and thin-factor ($0.85$)
  bars. (\textbf{b}) Per-patient coordinate agreement (Pearson $r$); all eight axes exceed $0.90$.
  The full loading scatter ($r=0.993$), the $\Phi$ comparison, and the model and inference details
  are in Supplement~\ref{ann:variational}.}
  \label{fig:vinuts}
\end{figure}
